% Part: first-order-logic
% Chapter: natural-deduction
% Section: quantifier-rules

\documentclass[../../../include/open-logic-section]{subfiles}

\begin{document}

\olfileid{fol}{ntd}{qrl}

\olsection{قواعد سورها}

\subsection{قواعد $\lforall$}

\begin{defish}
\AxiomC{$!A(a)$}
\RightLabel{\Intro{\lforall}}
\UnaryInfC{$\lforall[x][\Atom{!A}{x}]$}
\DisplayProof
\hfill
\AxiomC{$\lforall[x][\Atom{!A}{x}]$}
\RightLabel{\Elim{\lforall}}
\UnaryInfC{$!A(t)$}
\DisplayProof
\end{defish}

در قواعد~$\lforall$، $t$ ترمی بسته است (ترمی که هیچ متغیری ندارد) و
$a$~!!a{constant} است که نه در نتیجهٔ~$\lforall[x][!A(x)]$ رخ می‌دهد و
نه در هیچ فرضِ !!{undischarged} در !!{derivation} منتهی به
مقدمهٔ~$!A(a)$. $a$ را \emph{متغیر ویژهٔ} استنتاج
\Intro{\lforall} می‌نامیم.\footnote{اصطلاح «متغیر ویژه» را به کار
می‌بریم، هرچند $a$ در قاعدهٔ بالا یک ثابت است. این کاربرد دلایل تاریخی
دارد.}

\subsection{قواعد $\lexists$}

\begin{defish}
\AxiomC{$\Atom{!A}{t}$}
\RightLabel{\Intro{\lexists}}
\UnaryInfC{$\lexists[x][\Atom{!A}{x}]$}
\DisplayProof
\hfill
\AxiomC{$\lexists[x][\Atom{!A}{x}]$}
\AxiomC{[$\Atom{!A}{a}$]$^n$}
\DeduceC{$!C$}
\DischargeRule{\Elim{\lexists}}{n}
\BinaryInfC{$!C$}
\DisplayProof
\end{defish}

باز هم $t$ ترمی بسته است و $a$~!!a{constant} است که در مقدمهٔ
$\lexists[x][!A(x)]$، در نتیجهٔ~$!C$، یا در هیچ فرضِ
!!{undischarged} در !!{derivation}s منتهی به دو مقدمه رخ نمی‌دهد
(جز فرض‌های $!A(a)$). $a$ را \emph{متغیر ویژهٔ} استنتاج
\Elim{\lexists} می‌نامیم.

شرط اینکه متغیر ویژه نه در نتیجه‌ها و نه در هیچ فرضِ !!{undischarged}
در !!{derivation}s منتهی به مقدمات استنتاج \Intro{\lforall} یا
\Elim{\lexists} رخ دهد، \emph{شرط متغیر ویژه}
نامیده می‌شود.

\begin{explain}
قراردادمان را به یاد آورید که هرگاه $!A$~!!a{formula} باشد که
!!{variable}~$x$ در آن آزاد است، این امر را با نوشتن~$!A(x)$ نشان
می‌دهیم. در همان بافت، $!A(t)$ مخفف~$\Subst{!A}{t}{x}$ است. پس می‌توانیم
قاعدهٔ $\Intro\lexists$ را نیز چنین بنویسیم:
\begin{prooftree}
\AxiomC{$\Subst{!A}{t}{x}$}
\RightLabel{\Intro{\lexists}}
\UnaryInfC{$\lexists[x][!A]$}
\end{prooftree}
توجه کنید که ممکن است $t$ از پیش در~$!A$ رخ داده باشد؛ برای مثال، ممکن
است $!A$~برابر با~$\Atom{\Obj P}{t,x}$ باشد. بنابراین استنتاج
$\lexists[x][\Atom{\Obj P}{t,x}]$ از~$\Atom{\Obj P}{t,t}$ کاربردی درست
از~$\Intro\lexists$ است---می‌توانید یک یا چند رخداد~$t$ در مقدمه را،
و نه لزوماً همهٔ آن‌ها را، با !!{variable} مقید~$x$ «جایگزین» کنید. با
این حال، شرایط متغیر ویژه در $\Intro\lforall$ و~$\Elim\lexists$ ایجاب
می‌کنند که !!{constant}~$a$ در~$!A$ رخ ندهد. پس نمی‌توان
$\lforall[x][\Atom{\Obj P}{a,x}]$ را از $\Atom{\Obj P}{a,a}$ با
استفاده از~$\Intro\lforall$ به‌درستی استنتاج کرد.
\end{explain}

\begin{explain}
در \Intro{\lexists} و \Elim{\lforall} هیچ محدودیتی وجود ندارد و
ترم~$t$ می‌تواند هر چیزی باشد؛ پس لازم نیست نگران هیچ شرطی باشیم. از
سوی دیگر، در قواعد \Elim{\lexists} و \Intro{\lforall}، شرط متغیر ویژه
ایجاب می‌کند که !!{constant}~$a$ هیچ‌جا در نتیجه یا در یک فرض
!!{undischarged} رخ ندهد. این شرط برای تضمین صحت دستگاه ضروری است؛ یعنی
در آن فقط !!{derive}s کردنِ !!{sentence}s از فرض‌های !!{undischarged}ی
مجاز باشد که آن جمله‌ها از آن‌ها نتیجه می‌شوند. بدون این شرط، عبارت زیر
مجاز می‌بود:
\begin{prooftree}
  \AxiomC{$\lexists[x][!A(x)]$}
  \AxiomC{$\Discharge{!A(a)}{1}$}
  \RightLabel{*\Intro{\lforall}}
  \UnaryInfC{$\lforall[x][!A(x)]$}
  \RightLabel{\Elim{\lexists}}
  \BinaryInfC{$\lforall[x][!A(x)]$}
\end{prooftree}
اما $\lexists[x][!A(x)] \Entails/ \lforall[x][!A(x)]$.

از آنجا که قواعد حذف سورها تنها جایگذاری ترم‌های بسته به‌جای
!!{variable}s را مجاز می‌دانند، نتیجه می‌شود که هر !!{formula} که از
مجموعه‌ای از !!{sentence}s اشتقاق‌پذیر باشد خود !!a{sentence} است.
\end{explain}


\end{document}
